\footnote{A few of them nevertheless form stable compounds, for example with fluorine or oxygen. Such valences may arise when electrons are promoted out of the filled outer shell. The electrons then enter unoccupied $f$- (or $d$-) states lying relatively close in energy.

We should also note a distinctive attraction. It appears when a noble-gas atom interacts with an excited atom of the same element. Two identical atoms are then brought together while occupying different states. The number of possible states produced in this way is doubled (see \S~80). Transfer of the excitation from one atom to the other takes the place of the exchange interaction that produces ordinary valence. The molecule $\mathrm{He}_2$ provides an example. Bonding of the same kind occurs in molecular ions made up of two identical atoms. One example is $\mathrm H_2^+$.}
